\section{Abordagem Proposta}
\label{sec:metodo}

\subsection{Visão geral do pipeline}
\label{sec:metodo-visao}

\input{figures/architecture}

A abordagem estima como uma pessoa se posicionaria sobre um assunto numa audiência posterior, a
partir do que ela disse em audiências anteriores. Uma estimativa
desse tipo pode apoiar a preparação de audiências e ferramentas de transparência que informem o que
um participante já defendeu sobre um tema. Para conferi-la com o que a matéria atribuiu à pessoa, é
preciso saber que trecho da transcrição originou cada opinião, o que o PublicHearingBR não registra.
O pipeline (Figura~\ref{fig:arquitetura}) parte dos turnos de fala, do resumo estruturado e da
partição temporal das audiências (Seção~\ref{sec:dados-particao}).

A UDV (Seção~\ref{sec:metodo-udv}) associa o nome que a matéria cita aos turnos de fala da pessoa
naquela audiência e escolhe como evidência de cada opinião uma sentença dessa fala. É a sentença com
o trecho entre aspas da opinião ou, sem citação encontrada com segurança, a de sentido mais próximo, e
a sua localização na transcrição fica registrada para consulta.

Para escrever o perfil, o registro de ator (Seção~\ref{sec:metodo-atores}) reúne os turnos da mesma
pessoa em audiências diferentes, e um LLM recebe as falas dela nas audiências de treino. O perfil tem
até quatro blocos: posições, critérios e valores, alinhamentos declarados e forma de argumentar
(Seção~\ref{sec:metodo-perfis}). Como a UDV e o registro de ator localizam as falas pelas mesmas
referências (audiência, turno e posição na transcrição), o turno da evidência identifica o ator, e
dizemos que a opinião está ligada a ele.

Na simulação (Seção~\ref{sec:metodo-simulacao}), um LLM, o modelo simulador, responde como a pessoa
em três condições, que acrescentam o material um elemento de cada vez. A condição 0 recebe só o
nome e o cargo, a 1 acrescenta o perfil e a 2, trechos das falas de treino recuperados pelo assunto
da audiência.

Na avaliação (Seção~\ref{sec:setup}), cada opinião de validação ou de teste ligada a um ator com
perfil gera uma pergunta: entre a opinião e as de três outros participantes da mesma audiência, qual
a matéria atribuiu à pessoa. As perguntas de validação escolhem o número de trechos, e as de teste
medem o efeito de cada elemento; nenhuma fala dessas audiências entra no perfil ou nos trechos.

\subsection{Unidades Deliberativas Verificáveis}
\label{sec:metodo-udv}

As opiniões do PublicHearingBR registram o que a matéria atribui a cada pessoa, mas o conjunto não
indica que trecho da transcrição corresponde a cada opinião (Seção~\ref{sec:dados}). Sem essa
indicação, não é possível apresentar ao leitor a passagem da fala a que a opinião se refere nem
verificar se a fala a sustenta. A UDV estabelece o vínculo entre opinião e fala com um registro por
opinião que aponta, quando existe, o trecho da transcrição usado como evidência
(Figura~\ref{fig:udv}). A evidência está sempre contida num único turno da pessoa (uma sentença ou
sentenças contíguas), porque a segmentação em sentenças é feita dentro de cada turno e nenhuma
sentença junta falas separadas por falas de outros participantes. As regras completas de construção,
com os números de cada etapa, estão no Apêndice~\ref{ap:regras-udv}, e as ressalvas sobre a
resolução da pessoa, o corte de similaridade e o alcance do verificador, na
Seção~\ref{sec:limitacoes}.

\textbf{Resolução da pessoa.} O resumo estruturado registra o nome usado na matéria, que costuma
diferir do cabeçalho da transcrição (``Rodrigo Marinho'' na matéria, ``RODRIGO SARAIVA MARINHO'' na
transcrição), e quem preside a sessão aparece como \textit{PRESIDENTE}, com o nome entre parênteses.
A associação compara os conjuntos de palavras dos nomes, sem acento e em maiúsculas: um participante
recebe os turnos cujo nome é igual ao seu, contém o seu ou está contido nele, com tratamento próprio
para nomes de uma palavra. Com essa regra, 1.020 dos 1.065 participantes são associados a pelo menos
um turno, e os 45 restantes respondem por 90 opiniões.

\input{figures/udv-example}

\textbf{Citação literal.} A matéria apresenta como fala literal os trechos entre aspas das opiniões,
o que permite um vínculo por casamento de texto, sem modelo, mas pode editá-los. Na
Figura~\ref{fig:udv}, ela escreve ``presidente da comissão, que é eleita'', e a transcrição registra
``presidente que foi eleita''. Por isso, são buscados nos turnos da pessoa prefixos da citação, sem
exigir que o restante dela coincida com a fala, e prevalece o prefixo encontrado mais longo. Quando
esse prefixo tem seis palavras ou mais, a opinião tem uma \emph{citação confiável}: o nível é
\texttt{quote\_found}, e a evidência é a sentença que contém o prefixo.

Prefixos de três ou quatro palavras costumam ser expressões comuns, encontradas em pontos da fala sem
relação com a citação. O limite de seis palavras foi definido a partir de uma leitura exploratória
dos casamentos agrupados por tamanho de prefixo, sem protocolo de anotação. Das 2.113 opiniões de
pessoas associadas a turnos, 960 têm alguma citação com 10 caracteres ou mais; dessas, 549 têm algum
prefixo encontrado na fala da pessoa e 277 têm prefixo de seis palavras ou mais.

\textbf{Similaridade semântica.} Quando a opinião não tem citação confiável, ela e as sentenças
candidatas da pessoa são codificadas pelo Serafim 335m~\citep{gomes2024serafim}, um codificador de
sentenças para o português, e a evidência é a sentença de maior similaridade de cosseno com a
opinião. Um codificador denso foi preferido ao TF-IDF porque representa o sentido da sentença em
contexto, e a matéria costuma parafrasear a fala, caso em que a comparação por vocabulário
compartilhado falha.

\textbf{Nível e proveniência.} O nível resume o tipo de vínculo encontrado
(Tabela~\ref{tab:udv-niveis}), e os níveis \texttt{semantic\_match\_high} e
\texttt{semantic\_match\_weak} indicam de que lado do corte ficou o cosseno da melhor sentença, sem
afirmar que ela sustenta a opinião. A proveniência registra a origem do vínculo: \texttt{weak}
(rótulo produzido por regra, sem conferência humana) para o casamento de texto e \texttt{model} para
o codificador. Nenhum registro tem proveniência humana. O valor \texttt{weak} da proveniência não tem
relação com o nível \texttt{semantic\_match\_weak} nem com a corroboração fraca
(Apêndice~\ref{ap:regras-udv}).

\textbf{Localização e cobertura.} Todas as 2.105 evidências foram localizadas na transcrição
original, dentro do turno de origem, com offsets de caracteres, e as 98 opiniões restantes não têm
evidência. Nos 8 registros \texttt{no\_evidence}, a única fala transcrita da pessoa é a marcação
``(Manifestação em LIBRAS.)''. Entre as pessoas sem turno associado (\texttt{person\_not\_resolved}),
há casos de grafia divergente do nome, de instituições citadas como participantes e de pessoas que
não falaram na sessão. Nenhum desses níveis indica erro da matéria.

\input{tables/udv-tiers}

\textbf{Calibração do corte.} O corte de 0,45 (IC 95\% por \textit{bootstrap} de audiências, de 0,437
a 0,469) entre \texttt{semantic\_match\_high} e \texttt{semantic\_match\_weak} foi calibrado nas
audiências de treino, porque a amostra de validação humana (Seção~\ref{sec:setup-udv}) é sorteada
no teste e depende dessa fronteira. Ele vem de uma regra de pares, com positivos formados por cada
uma das 183 opiniões de treino com citação confiável e a sentença da própria citação, e negativos
formados pela mesma opinião e uma sentença sorteada de outra audiência do treino. Essa regra, que toma
o ponto médio entre o primeiro quartil dos cossenos dos positivos e o terceiro quartil dos negativos,
foi adotada depois que a regra declarada como principal, pelo índice J de
Youden~\citep{youden1950index}, se mostrou degenerada (Apêndice~\ref{ap:calibracao}). Como 43 das
1.828 UDVs semânticas (2,4\%) ficam abaixo do corte, o nível \texttt{semantic\_match\_weak} reúne
poucos registros.

\textbf{Verificação.} Um verificador separado refaz, a partir das transcrições e dos resumos
estruturados, a resolução de pessoa, a extração das sentenças candidatas e a busca de citações.
Em seguida, confere se a opinião e a pessoa coincidem com o resumo estruturado, se a evidência está
num único turno da pessoa, se os offsets recuperam o texto e se o nível é coerente com a busca de
citações, o cosseno e o corte. Como reutiliza os procedimentos da construção nessas três etapas, o
verificador confere a consistência dos registros com esses procedimentos, e não a correção deles
(Apêndice~\ref{ap:regras-udv}); no conjunto descrito aqui, não acusou nenhuma inconsistência.

\subsection{Falas por ator entre audiências}
\label{sec:metodo-atores}

Como a UDV trata uma audiência de cada vez, descrever o que uma pessoa defende ao longo do tempo
exige reunir suas falas em todas as audiências de que participou e, para isso, decidir quando
cabeçalhos de audiências diferentes se referem à mesma pessoa. A regra de nome contido da UDV compara
nomes dentro de uma única audiência, que tem poucos falantes; aplicada entre audiências, ela uniria
pessoas diferentes, como ``Luiz Lima'' e ``Gustavo Luiz de Lima Correia''. Por isso, o registro de
ator usa como chave o nome do cabeçalho sem acento e em maiúsculas (nos turnos de presidência, o nome
entre parênteses), complementado por uma lista de mesclas.

A lista de mesclas foi proposta a partir dos 153 pares de nomes candidatos, em que um nome contém o
outro ou os dois têm as mesmas palavras, e validada com evidências do próprio conjunto
(Apêndice~\ref{ap:regras-atores}). O resultado tem 37 grupos de mescla, uma chave dividida por
audiência (um caso de homônimo exato) e um par sem resolução, mantido separado.

São removidos os turnos de intérprete e de tradução, cuja fala pertence a outro participante, os
turnos vazios ou formados por marcação de palco e os turnos de presidência com menos de 50 palavras,
tratados como condução da sessão. Os turnos de presidência concentram fala de condução, mas
descartá-los por completo eliminaria opiniões, pois 336 das 2.105 evidências de UDV estão nesses
turnos, e o limite de 50 palavras mantém 332 delas. Após os filtros, restam 13.190 dos 17.264 turnos
(Apêndice~\ref{ap:regras-atores}).

Os 301 atores que falam em duas ou mais audiências formam o material dos perfis, porque neles é
possível escrever o perfil com parte das audiências e avaliá-lo nas demais. Desses, 264 falam em
pelo menos uma audiência de treino e recebem perfil (Seção~\ref{sec:dados-particao}), 79 deles com
uma única audiência de treino.

\subsection{Perfis por ator}
\label{sec:metodo-perfis}

A fala de um ator se distribui por audiências de temas diferentes e se mistura a passagens de
condução da sessão e de cortesia. O perfil condensa essa fala num texto escrito por um LLM, que serve
de instrução ao modelo simulador (Seção~\ref{sec:metodo-simulacao}), tanto sobre temas que a
pessoa já tratou quanto sobre temas novos. Por isso, além do que a pessoa defende, critica, propõe e
cobra, o perfil registra os critérios que ela invoca para justificar essas posições e o que declara
sobre si e sobre o próprio lado. O perfil registra também a forma como a pessoa argumenta
(Figura~\ref{fig:perfil}). Cada perfil usado na avaliação é escrito apenas com as audiências de
treino (Seção~\ref{sec:dados-particao}), para que nenhuma fala ou opinião das audiências de validação
e de teste chegue ao prompt.

\input{figures/profile-example}

Cada ator recebe uma única geração, cujo \textit{user prompt} contém todas as suas falas de treino,
agrupadas por audiência e em ordem cronológica. As falas de cada audiência vêm depois de uma linha
com a data da matéria e o assunto registrado no resumo estruturado, e os turnos de presidência são
precedidos da marca \texttt{[presidência da sessão]}. A data permite situar no tempo eventuais
mudanças de posição, e o assunto ajuda a resolver referências vagas, como ``este projeto de lei''. O
\textit{user prompt} identifica a pessoa apenas pelo nome, sem o partido e a UF dos cabeçalhos dos
seus turnos, e o partido só chega ao LLM quando a própria pessoa o menciona nas falas. O
\textit{system prompt} impõe quatro grupos de regras, e os dois prompts estão no
Apêndice~\ref{ap:prompt}.

As regras de fonte única restringem o perfil às falas, sem conhecimento prévio sobre a pessoa, os
partidos ou os temas em debate. Como o assunto vem do resumo estruturado, e não da fala, nada que
apareça só no assunto é atribuído à pessoa. Propostas de terceiros que ela relata ou questiona não se
tornam posições dela, e siglas, nomes e datas não são completados. As regras de condução e cortesia
determinam que sejam desconsideradas passagens como passar a palavra, controlar o tempo e agradecer.

As regras de blocos organizam o perfil em até quatro blocos (posições; critérios e valores;
alinhamentos declarados; forma de argumentar), omitem o bloco sem apoio nas falas e proíbem deduzir
valores ou alinhamentos das posições ou do partido. Na forma de argumentar, um traço precisa aparecer
em mais de um turno sem ser contrariado. As regras de escrita pedem itens ordenados por importância,
proíbem registrar ausências, como o que a pessoa não mencionou ou deixou de defender, e fixam um teto
de 800 palavras.

A divisão em blocos destina-se a separar a parte do perfil que orienta respostas sobre temas já
tratados (as posições) da parte que orienta respostas sobre temas novos (os critérios e os
alinhamentos declarados), aos quais o \textit{system prompt} da simulação remete o modelo quando o
material da conversa não trata do assunto (Seção~\ref{sec:metodo-simulacao}). A forma de argumentar
também chega ao modelo simulador como parte do perfil, mas nenhuma instrução da simulação se refere a
ela. As proibições de inferência visam impedir que o perfil atribua à pessoa características do partido ou das posições
que ela própria não enunciou. Com um alinhamento deduzido do partido, o modelo que recebe o perfil
responderia com base nele, e não nas falas. A proibição de registrar ausências tem o mesmo motivo,
porque uma frase como ``não menciona temas de saúde'' poderia ser tomada como posição.

As regras sobre o assunto da audiência, sobre posições de terceiros, sobre siglas e nomes e sobre o
bloco de forma de argumentar foram acrescentadas após a leitura de perfis gerados durante o desenvolvimento do
prompt. Esses perfis foram gerados por outros LLMs e, em parte, com falas de todas as audiências,
inclusive as de teste. Neles, o assunto da audiência aparecia como posição do ator, a proposta de um
convidado resumida pelo presidente aparecia como proposta do próprio presidente, e siglas e sobrenomes
eram completados com conhecimento externo às falas. Além disso, uma frase dirigida a uma única
autoridade aparecia como traço geral da pessoa, e um traço atribuído a ela era contrariado por falas
de outra audiência. As regras resultantes tratam da forma do perfil, e não do conteúdo de posições
específicas, e os perfis de desenvolvimento não entram na avaliação.

\subsection{Simulação de atores}
\label{sec:metodo-simulacao}

Na simulação, o perfil é fornecido a um LLM, o modelo simulador, que deve responder como a pessoa
responderia; nos experimentos, cada um dos dois LLMs avaliados simula com os perfis que ele próprio
escreveu (Seção~\ref{sec:setup-simulacao}). O simulador é avaliado
por perguntas de múltipla escolha, em que escolhe entre opiniões publicadas sobre a audiência
(Seção~\ref{sec:setup-simulacao} e Figura~\ref{fig:simulacao}). Dois elementos, o perfil e os trechos recuperados, são
acrescentados um de cada vez, o que define as três condições encadeadas da
Tabela~\ref{tab:condicoes}. A condição 0 recebe o nome e o cargo da pessoa, sem outro material; a 1
acrescenta o perfil; e a 2, trechos das falas de treino recuperados pelo assunto.
% Pendente: motivo de usar o mesmo LLM para gerar os perfis e para simular.

\input{tables/conditions}

\textbf{Prompt comum.} Todas as chamadas partem do mesmo \textit{system prompt} e diferem no material
incluído (Apêndice~\ref{ap:prompt-simulacao}). Com instruções diferentes por condição, a diferença
entre condições vizinhas mediria também a troca de instrução. O prompt apresenta a pessoa pelo nome
e pelo cargo registrado na matéria da audiência avaliada. As instruções proíbem o modelo de usar
conhecimento prévio sobre a pessoa, os partidos e o grupo que ela representa, assim como a própria
opinião. Quando o material não trata do assunto, as instruções orientam o modelo a se basear nos
critérios e nos alinhamentos da pessoa.

O \textit{system prompt} também instrui o modelo a escolher, entre as opções dadas, a mais compatível com
o material, mesmo que nenhuma apareça nele. Essa instrução foi acrescentada, e o texto do pedido de
múltipla escolha (Seção~\ref{sec:setup-simulacao}) reescrito, depois de um teste preliminar com dois
atores e três modelos, que incluiu perguntas de audiências de validação e de teste.
% Pendente: declarar se os dois atores do teste preliminar entram na avaliação final.

O partido não é fornecido em nenhuma condição, para que a falta de material não seja suprida pela
posição associada a ele. A sigla entre parênteses, com ou sem a UF, é retirada do cargo (``Deputado
(PT-SP)'' torna-se ``Deputado''), e o partido só chega ao modelo quando a própria pessoa o declara
nas falas. Cada item do perfil no formato pedido recebe um identificador formado pela inicial do
bloco e pela ordem do item (P1, C1, A1, F1). Na condição 0, o perfil é substituído por ``Nenhum.''.

\textbf{Trechos recuperados.} Como o perfil condensa as falas de treino num texto com teto de 800
palavras, detalhes que podem decidir uma pergunta são omitidos. Para recuperá-los, o assunto da
audiência é codificado pelo Serafim 335m da UDV e comparado às sentenças de quatro palavras ou mais
dos turnos da pessoa nas audiências de treino. São incluídas as $k$ sentenças de maior cosseno, no
máximo uma por turno, cada uma com a sentença anterior e a seguinte do mesmo turno. Os trechos são
incluídos no \textit{user prompt} em ordem de data, para que uma mudança de posição apareça na ordem
em que ocorreu. Cada um traz um identificador (T1 a T$k$), a data e o assunto da audiência de origem
e, quando for o caso, a marca de presidência; o cosseno não é mostrado ao modelo.
